\begin{lemma}\label{lem:GrobnerLink}
With notation as above, if $w \in S_n$ avoids obstructions of Types $1$, $2$, and $3$ and Conjecture \ref{conj:mainresult} holds for all permutations of smaller Coxeter length than that of $w$, then either $D_w = \Dom(w)$ or there is some lower outside corner $(i,j)$ of $D_w \setminus \Dom(w)$ corresponding to the variable $y = z_{i,j}$ so that the generators $\{q_1, \ldots, q_k, h_1, \ldots, h_\ell\}$ of $C_{y,I_w}$ form a Gr\"obner basis under any diagonal term order.
\end{lemma}
\begin{proof}
Fix a diagonal term order $\sigma$, and assume $D_w \neq \Dom(w)$.  By Lemma \ref{lem:GrobnerQs}, there is some lower outside corner of $D_w$ so that, with notation as above, there exists $\ell' \geq 0$ so that the generators $\{q_1, \ldots, q_k, h_1, \ldots, h_{\ell'} \}$ form a Gr\"obner basis for $Q_{y,I_w}$.  

By Corollary \ref{cor:GrobnerDel} and the inductive hypothesis, $\{h_1,\ldots, h_\ell\}$ is a diagonal Gr\"obner basis for $N_{y,I_w}$.  Then by Lemma \ref{lem:concatinate}, the generators $\{q_1, \ldots, q_k, h_1, \ldots, h_\ell\}$ form a diagonal Gr\"obner basis for $C_{y,I_w}$ if and only if for every $q_a$ and $h_b$ there exists some $f \in (q_1, \ldots, q_k, h_1, \ldots, h_{\ell'}) \cap (h_1, \ldots, h_\ell)$ satisfying $LT(f) = LCM(LT(q_a),LT(h_b))$.  

If $h_b \in (q_1, \ldots, q_k, h_1, \ldots, h_{\ell'})$ or $q_a \in (h_1, \ldots, h_\ell)$, the result follows from Lemma \ref{lem:concatinate}, and if $LCM(LT(q_a),LT(h_b)) = LT(q_a) \cdot LT(h_b)$, then we take $f = q_a \cdot h_b$.  Otherwise, there is some $(r,s) \in \Ess(w) \setminus \Dom(m)$ with $\rank_w(r,s) = \deg(h_b)-1$ corresponding to the variable $e = z_{r,s}$ weakly southeast of all of the variables involved in $h_b$.  Because $w$ has no Type $3$ obstruction, $e$ is not strictly northwest of $y$.  Suppose first that $y$ is strictly east and weakly north of $e$.

In this case, let $M'$ be the matrix consisting of the union of the columns determining $q_a$ and $h_b$ and the union of the rows determining $q_a$ and $h_b$.  We will next describe an auxiliary matrix $M$ formed from $M'$. (For an example of the construction, see Example \ref{ex:formM} below.)  Form a matrix $M$ from $M'$ as follows: First set to $0$ any entry whose row index is not one of the rows determining $q_a$ and whose column index is not one of the columns determining $h_b$.  Next, whenever a column of $M'$ contains a variable dividing the leading term of $q_a$ and a distinct variable dividing the leading term of $h_b$, duplicate that column and replace in one copy of the column the variables coming only from $h_b$ by $0$.  Whenever a row of $M'$ contains a variable dividing the leading term of $q_a$ and a distinct variable dividing the leading term of $h_b$, duplicate that row and replace in one copy of the row the variables coming only from $q_a$ by $0$.  

Now $M$ will be a $d \times d$ matrix where $d = \deg(LCM(LT(q_a),LT(h_b)))$ because it will have one row and one column for each monomial dividing $LT(q_a)$ and one each for every monomial dividing $LT(h_b)$ but not $LT(q_a)$.  

By expressing $\det(M)$ as a sum of products of the $\deg(q_a)$-minors from the rows of $M$ originating from the submatrix of $Z_w$ determining $q_a$ and the $(d-\deg(q_a))$-minors in the remaining rows, we see that $\det(M) \in (q_1, \ldots, q_k)$.  Similarly, by expressing $\det(M)$ as a sum of products of $\deg(h_b)$-minors in the column originating from the submatrix of $Z_w$ determining from $h_b$ and $(d-\deg(h_b))$-minors in the remaining columns, we have $\det(M) \in (h_1, \ldots, h_\ell)$.  It is because $y$ is strictly east and weakly north of $e$ that the rows determining $q_a$ that every $\deg(q_a)$-minor in the specified rows is an element of $(q_1, \ldots, q_k)$ and that every $\deg(h_b)$-minor in the specified column is an element of $(h_1, \ldots, h_\ell)$.

Next we will see that $LT(\det(M)) = LCM(LT(q_a),LT(h_b))$.  Call $\tilde{M}$ the submatrix of $M'$ whose entries are northwest of both $e$ and $y$.  Set $\mu_1$ to be the product of the terms of $\tilde{M}$ that divide either $LT(q_a)$ or $LT(h_b)$.  Call $\mu_2$ the leading term of the determinant of the submatrix of $M$ consisting of the rows and columns used to determine $h_b$ excluding those involving a divisor of $\mu_1$.  Similarly, define $\mu_3$ to be the leading term of the determinant of the submatrix of $M$ consisting of the rows and columns determining $q_a$ excluding those involving a divisor of $\mu_1$.  

Notice that $\mu_1 \cdot \mu_2 \cdot \mu_3 = LCM(LT(q_a),LT(h_b))$ and that this product is a term of $\det(M)$.  To see that every other term of $\det(M)$ is smaller under $\sigma$, notice that because $w$ has no obstruction of Type $1$, the submatrix of $M'$ whose entries are both northwest of $e$ and northwest of $y$ must have $0$'s only in full rows and full columns along the north and west sides.  

It will now be more convenient for us to work with $\widehat{M}$, obtained from $M$ by adding the doubled copies of rows and columns obtained in the transition from $M'$ to $M$ so that no variable appears more than once.  Note that $\det(\widehat{M}) = \det(M)$.  If some other term of $\det(\widehat{M})$ is larger than $\mu_1 \cdot \mu_2 \cdot \mu_3$, there must be some entries of $\widehat{M}$ dividing $\mu_1 \cdot \mu_2 \cdot \mu_3$ whose row indices we may permute to obtain a larger monomial.  We may assume that this permutation consists of one cycle.  If all entries divide either $\mu_1 \cdot \mu_2$ or $\mu_1\cdot \mu_3$, we would obtain a term of $q_a$ or $h_b$, respectively, that is strictly larger than its leading term, which also cannot be.  But the permutation cannot send any divisor of $\mu_3$ to the row of a divisor of $\mu_1$, all of which are $0$ in that column, or vice versa.  Hence, $\mu_1 \cdot \mu_2 \cdot \mu_3 = LT(\det(\widehat{M})) = LT(\det(M))$, as desired.

Finally, if $y$ is strictly south and weakly west of $e$, a parallel argument gives the~result.
\end{proof}
